% KaTeX's alphabet selectors retain ordinary symbols outside their Latin repertoire.
\DeclareMathAlphabet{\mathsfbf}{OT1}{cmss}{bx}{n}
\DeclareMathAlphabet{\mathsfbfit}{OT1}{lmss}{bx}{sl}
% Unicode bold italic specifies both properties, independently of an enclosing alphabet.
\newcommand{\ReferenceBoldItalic}[1]{\bm{\mathnormal{#1}}}
\newcommand{\ReferenceBoldScript}[1]{\bm{\mathscr{#1}}}
\newcommand{\ReferenceBoldFraktur}[1]{\bm{\mathfrak{#1}}}
\newcommand{\ReferenceUnicode}[4]{\text{\fontfamily{#1}\fontseries{#2}\fontshape{#3}\selectfont #4}}
\makeatletter
\newcommand{\ReferenceFixedSymbol}[1]{%
  \count@=#1\relax
  \ifnum\count@>"6FFF \advance\count@ by -"7000 \mathchardef#1=\count@\fi}
\newcommand{\ReferenceKeepSymbols}{%
  \forcsvlist{\ReferenceFixedSymbol}{\Gamma,\Delta,\Theta,\Lambda,\Xi,\Pi,\Sigma,\Upsilon,\Phi,\Psi,\Omega,\imath,\jmath}}
\newcount\ReferenceCode
\newcommand{\ReferenceKeepRange}[3]{%
  \ReferenceCode=#1\relax
  \loop\mathcode\ReferenceCode=\numexpr#3*256+\ReferenceCode\relax
    \ifnum\ReferenceCode<#2\relax\advance\ReferenceCode by 1\repeat}
\let\ReferenceFraktur\mathfrak
\renewcommand{\mathfrak}[1]{{\ReferenceKeepSymbols\ReferenceFraktur{#1}}}
\let\ReferenceScript\mathscr
\renewcommand{\mathscr}[1]{{\ReferenceKeepSymbols
  \ReferenceKeepRange{`a}{`z}{1}\ReferenceKeepRange{`0}{`9}{0}\ReferenceScript{#1}}}
\makeatother

% mu is one eighteenth of the current math style's symbol-font quad.
\newdimen\ReferenceRaise
\newdimen\ReferenceWidth
\newdimen\ReferenceHeight
\ExplSyntaxOn
\cs_new_protected:Npn \reference_dimension_set:Nnn #1#2#3 {
  \regex_extract_once:nnNTF {\A\s*([+\-]?(?:\d+\.?\d*|\.\d+))mu\s*\Z} {#3} \l_tmpa_seq
    {
      \tl_set:Nx \l_tmpa_tl {\seq_item:Nn \l_tmpa_seq {2}}
      \dim_set:Nn #1 {\l_tmpa_tl\fontdimen6
        \ifx#2\scriptstyle\scriptfont\else
          \ifx#2\scriptscriptstyle\scriptscriptfont\else\textfont\fi\fi 2/18}
    }
    {\dim_set:Nn #1 {#3}}
}
\cs_new_eq:NN \ReferenceSetDimension \reference_dimension_set:Nnn
\ExplSyntaxOff
\newcommand{\ReferenceRule}[3][0pt]{%
  \ifmmode\mathpalette{\ReferenceRuleAux{#1}{#2}{#3}}{}%
  \else\rule[#1]{#2}{#3}\fi}
\newcommand{\ReferenceRuleAux}[5]{\begingroup
  \ReferenceSetDimension\ReferenceRaise{#4}{#1}%
  \ReferenceSetDimension\ReferenceWidth{#4}{#2}%
  \ReferenceSetDimension\ReferenceHeight{#4}{#3}%
  \rule[\ReferenceRaise]{\ReferenceWidth}{\ReferenceHeight}\endgroup}
